% ch-a4-16.tex —— 第 4 篇（Emerton--Gee--Hellmann）续片，底本 PDF 360--371 = 印刷页 347--358
\ysrc{347}

\begin{thm}{9.4.2（预期）}\label{thm:a4-16-9-4-2}
\textit{（M.E.--T.G.--Zhu）}设 $L_\infty$ 是预期定理 7.2.1 的 $\mathrm{PGL}_2(\Q_p)$-表示层。则存在 $\mathrm{PGL}_2(\Q_p)$-等变同构
\begin{equation}\tag{9.4.3}\label{eq:a4-16-9-4-3}
\wt H_{1,\frakm}\isom R\Gamma(\XX_{\bar r},f^{!}L_\infty\otimes\mathcal V)[-1],
\end{equation}
它与两侧 $\Gal_\Q$ 以及 $\ell\neq p$ 处（导出）Hecke 代数的作用相容。

此外，还有 $\mathrm{PGL}_2(\Q_p)$-等变同构 $M_\infty\isom g_*L_\infty$。
\end{thm}

\begin{rem}{9.4.4}\label{rem:a4-16-9-4-4}
我们预计，预期定理 9.4.2 中的假设可以大幅放宽。特别地，可以允许 $\bar r$ 在不整除 $p$ 的位置分歧；同时希望在 [EGZ] 定稿时尽可能放宽对 $\bar r$ 的其他假设。
\end{rem}

我们打算如下证明预期定理 9.4.2。基本想法是：$M_\infty$ 由模曲线的上同调「拼接」而成，而我们想说明，拼接过程会「抵消」上同调猜想公式中出现的 $f^{!}$。我们（利用已有的模性提升定理）证明 $\XX_{\bar r}$ 等于 $[\Spf R_{\{p,\infty\}}(\bar r)/\boldsymbol{\mu}_2]$。接着证明 $R_{\{p,\infty\}}(\bar r)$ 是 $R(\bar\rho)$ 对某个正则序列的商，而该序列同时也是 $M_\infty$ 上的正则序列；这样便可借助 Koszul 复形计算拉回 $f^{!}$，最终把 (9.4.3) 归结为 M.E. 最初在 [Eme11a] 中证明的局部-整体相容性结果，或者更确切地说，归结为 [CEG$^+$18] 对这些结果的解释与重新证明。

\subsection{$\XX^{\mathrm{odd}}_{\mathrm{GL}_2,\Q,S}$ 的结构}

[EGZ] 的另一个目标，是借助 Taylor--Wiles 拼接与 Iwasawa 理论，描述叠 $\XX^{\mathrm{odd}}_{\mathrm{GL}_2,\Q,S}$ 的结构。

我们简要描述若干预期结果。为简单起见，再次只限于讨论行列式为 $\varepsilon^{-1}$ 的叠 $\XX^{\varepsilon^{-1},\mathrm{odd}}_{\mathrm{GL}_2,\Q,S}$。

类似于 (9.1.1)，有分解
\[
\XX^{\mathrm{odd}}_{\mathrm{GL}_2,\Q,S}:=\coprod_{\bar\theta}\XX_{\bar\theta}
\]
为连通分支之并，其中 $\bar\theta$ 遍历 $\Gal_{\Q,S}$ 的、行列式为 $\bar\varepsilon^{-1}$ 的奇剩余伪表示（odd residual pseudorepresentation）。最简单的情形是 $\bar\theta$ 来自不可约奇表示 $\bar r:\Gal_{\Q,S}\to\mathrm{GL}_2(\F_p)$。此时记 $X_{\bar\theta}=X_{\bar r}=\Spf R_{\bar r}$，其中 $R_{\bar r}$ 是 $\bar r$ 的形变环；并发现 $\XX_{\bar\theta}=\XX_{\bar r}=[\Spf R_{\bar r}/\wh{\GG}_m]$，其中 $\wh{\GG}_m$ 表示 $\GG_m$ 的 $\frakm_{R_{\bar r}}$-adic 完备化，它平凡地作用在 $\Spf R_{\bar r}$ 上。于是在此情形，模叠 $\XX_{\bar r}$ 的理论与通常的伽罗瓦形变论（Galois deformation theory）没有本质差别。Serre 猜想告诉我们 $\bar r$ 是模的；进而已知的「big $R$ = big $\mathbf T$」型定理蕴含（至少在 $\bar r$ 满足通常的 Taylor--Wiles 假设时）：$R_{\bar r}$ 可以等同于某个适当定义的 Hecke 代数，该代数据作用在 $p$-进模形式的空间上。

当 $\bar\theta$ 可约时，情形更为微妙；此时的目标仍是在模形式的意义下尽可能描述 $\XX_{\bar\theta}$。为说明其含义，令 $\bar\theta=1+\bar\varepsilon^{-1}$，并令 $\mathcal Y_{\bar\theta}$ 为（导出）子叠，它由如下局部条件刻画：在素数 $\ell\in S\setminus\{p,\infty\}$ 处幂幺驯顺，且在 $p$ 处有限平坦。令 $N=\prod\limits_{\ell\in S\setminus\{p,\infty\}}\ell$，并选取足够大的 $\OO$，使得每个权为 $2$、级为 $N$ 的尖点特征形式（cuspidal eigenform）$f$ 的系数都落在 $\OO$ 中。于是每个这样的尖点形式 $f$，只要其相伴的模 $p$ 伪表示 $\theta_f$ 等于 $1+\bar\varepsilon^{-1}$，就贡献 $\mathcal Y_{\bar\theta}$ 的一个不可约分支，其形状为 $[\Spf\OO\langle X,Y\rangle/(XY-c)/\wh{\GG}_m]$，其中 $c$ 是某个可约化模（modulus of reducibility），$\OO\langle X,Y\rangle$ 表示 $\OO[X,Y]$ 的 $p$-adic 完备化，而 $\wh{\GG}_m$ 表示 $\GG_m$ 的 $p$-adic 完备化，其作用为 $t\cdot(X,Y)=(t^2X,t^{-2}Y)$。这个分支参数化与 $f$ 相伴的 $p$-adic 伽罗瓦表示中的格。

\ysrc{348}
除这类分支外，还有一个「Eisenstein」分支，其底层的经典叠形如 $[\Spf\OO\langle X_1,\dots,X_n\rangle/\wh{\GG}_m]$，其中 $n$ 表示 $S\setminus\{p,\infty\}$ 的阶（即 $N$ 的因子个数），而 $\wh{\GG}_m$ 的作用为 $t\cdot X_i=t^2X_i$。

这个分支参数化 $\varepsilon^{-1}$ 被 $1$ 扩张的情形；这类扩张的空间由与 $N$ 的 $n$ 个因子相伴的 Kummer 类张成。若 $n>1$，则该分支的维数大于预期，于是（按上面的讨论）会带有非平凡的导出结构。

$\XX$ 还可能有另一个不可约分支，它参数化 $1$ 被 $\varepsilon^{-1}$ 扩张的情形（即「方向反了」的扩张）。为简单起见设 $p\ge5$，这类扩张由那些满足 $\ell\equiv-1\bmod p$ 的 $\ell\in S$ 控制。若有 $m$ 个这样的 $\ell$，则这类扩张构成的族是 $(m-1)$ 维的（特征 $p$ 中有 $m$ 维的扩张空间，再对这个空间取 $\wh{\GG}_m$-作用的商），并且其支撑（作为集合）完全位于特征 $p$ 中。若 $m>1$，它就是 $\mathcal Y_{\bar\theta}$ 的一个分支；若 $m>2$，它的维数同样大于预期，因而容许导出结构。

我们预期上述讨论穷尽了 $\mathcal Y_{\bar\theta}$ 的所有可能分支。我们希望，通过适当应用 Fontaine--Mazur 猜想（的已知情形，特别地，剩余可约情形已由 Skinner--Wiles [SW97] 与 Pan [Pan22a] 证明），并借助 Iwasawa 理论的主猜想 [MW84]，可以证明这一点（以及更一般的陈述）。

\subsection{特征簇与超收敛 $p$-adic 自守形式的层}

我们现在陈述猜想 9.3.2 的一个类比，对某些（超收敛）$p$-adic 系数系统的上同调给出猜想性的描述。由此便得到特征簇（eigenvariety）上出现的超收敛 $p$-adic 自守形式之层的猜想性描述。我们沿用前一节的记号。

处理 $p$-adic 自守形式与特征簇基本上有两条途径：其一是使用局部对称空间上、系数取 $p$-adic 系数系统的上同调（例如见 [Che04] 或 [Han17]）；其二是使用 M.E. 引入、并由 M.E. 与 Frank Calegari 进一步发展的完备上同调群（见 [Eme06b, Eme14, CE12]），并施用 M.E. 在 [Eme06a] 中引入的局部解析 Jacquet 函子（用这条途径构造特征簇见 [Eme06b]）。下面简要讨论这两条途径之间的关系，以及它们与猜想 6.2.4 中函子 $\mathfrak A_G^{\mathrm{rig}}$ 的关系。

\begin{rem}{9.6.1}\label{rem:a4-16-9-6-1}
一般地，特征簇可以与如下数据相伴：驯顺级 $K^p$，$v|p$ 处 $\Rep_L^{\mathrm{sm}}G(F_v)$ 的 Bernstein 块 $\Omega_v=[M_v,\sigma_v]$，以及 $\Omega_v$ 的超尖点支撑 $M_v$ 的代数表示 $W_v$（其定义只差 $M_v$ 的一个代数特征标之扭）。我们将聚焦于经典情形，即所有 Bernstein 块都是主的（此时 $W_v$ 不是额外数据），至于一般 Bernstein 分量的处理，请参看 [BD24]。我们指出，在这一更一般的背景下，也仍可陈述与下文猜想 9.6.9 和猜想 9.6.31 类似的猜想。
\end{rem}

\ysrc{349}
\subsubsection*{9.6.2.\quad 志村集情形}

我们先讨论特殊情形：约化群 $G$ 在无穷位紧，从而局部对称空间 $Y(K_f)$ 只是点集。我们把这类集合称为志村集（Shimura set），尽管严格说来它们并非志村簇。为固定设定，从今往后假设 $F$ 全实，且 $G$ 是
\begin{itemize}
\item 与 $F$ 的某个 CM 扩张 $E$ 上的 Hermite 空间相伴的定酉群（definite unitary group），
\item 或与 $F$ 上某个定四元数代数相伴的群；
\end{itemize}
并且 $G$ 在整除 $p$ 的位置分裂。

对每个 $v|p$ 的位置，取定一个分裂极大环面以及一个 Borel 子群 $T_v\subset B_v\subset G_v=G(F_v)=\mathrm{GL}_d(F_v)$。我们还取定适应 $B_v$ 的 Iwahori 子群 $I_v$，并记 $T_v^0=T_v\cap I_v$ 为环面 $T_v$ 的极大紧子群，$B_{v,0}=B_v\cap I_v$。此外，取定驯顺级 $K^p$，它是 $G(\mathbb{A}_F^{\infty,p})$ 的紧开子群。设 $\kappa_v:T_v^0\to L^\times$ 是局部解析特征标。于是可以定义如下取值于模的超收敛 $p$-adic 自守形式空间
\begin{equation}\tag{9.6.3}\label{eq:a4-16-9-6-3}
\begin{aligned}
M^\dagger(&G,K^p,(\kappa_v)_{v|p})=\Bigl\{f:G(F)\backslash G(\mathbb{A}_F^\infty)/K^p\longrightarrow\bigotimes_{v|p}\bigl(\Ind_{B_{v,0}}^{I_v}\kappa_v\bigr)^{\mathrm{an}}\ \Bigm|\\
&\text{满足 } f(gu)=u^{-1}f(g)\ \text{对一切 } g\in G(\mathbb{A}_F^\infty)\text{ 与 } u\in\prod_{v|p}I_v\Bigr\},
\end{aligned}
\end{equation}
可对照 [Loe11, 3.3]。该空间带有所谓 Atkin--Lehner 环 $\mathcal A^+(G_v)$ 的作用，下面简述其定义（及作用）。

对每个位置 $v|p$，Borel 子群 $B_v$ 及其子群 $N_{v,0}=N_v\cap B_{v,0}$ 的选取决定了一个子幺半群 $T_v^+\subset T_v$，它由所有满足 $tN_{v,0}t^{-1}\subset N_{v,0}$ 的 $t\in T_v$ 组成。显然 $T_v^+$ 包含极大紧子群 $T_v^0$。取定自由阿贝尔子群 $\Sigma_v\subset T_v$，使得 $\Sigma_v\cong T_v/T_v^0$（在对角环面情形，选定单值化元 $\varpi_v$ 之后，例如可以取所有元素均为 $\varpi_v$ 的幂的对角矩阵），并令 $\Sigma_v^+=\Sigma_v\cap T_v^+$，它是自由阿贝尔幺半群。然后定义
\[
\mathcal A^+(G_v)=L[\Sigma_v^+],\qquad
\mathcal A^+(G)=\bigotimes_{v|p}\mathcal A^+(G_v).
\]
由 [Loe11, 3.4] 可知，$\mathcal A^+(G_v)$ 典范同构于 Hecke 代数 $\mathcal H(G_v)$ 中由 $I_v$-双不变、且支撑在 $G_v$ 的由 $\Sigma_v^+$ 与 $I_v$ 生成的子幺半群 $\Sigma_v^+I_v$ 上的函数所构成的子代数。$I_v$ 在 $\Ind_{B_{v,0}}^{I_v}\kappa_v$ 上的作用典范地延拓到 $\Sigma_v^+I_v$ [Loe11, Theorem 2.4.7]，于是该幺半群通过 $(\gamma.f)(g)=\gamma f(g\gamma)$ 作用在 $M^\dagger(G,K^p,(\kappa_v)_{v|p})$ 上。把 $\mathcal A^+(G)$ 的元素看作 $G(\prod_{v|p}F_v)$ 上的局部常值函数，便得到 $\mathcal A^+(G)$ 在 $M^\dagger(G,K^p,(\kappa_v)_{v|p})$ 上的作用。

下面给出上述超收敛 $p$-adic 自守形式空间及其 $\mathcal A^+(G)$-作用的另一种刻画；这种刻画涉及完备上同调，写起来其实更省事（尤其在这个框架中，Atkin--Lehner 环的作用更容易写出）。不过我们仍先给出上面的定义，因为它与系数取

\ysrc{350}由表示 $\Ind_{B_{v,0}}^{I_v}\kappa_v$ 定义的 $p$-adic 系数系统的上同调直接相关（这使得下文猜想 9.6.9 的表述更容易）。

考虑空间
\[
\wh S(K^p,L)=\{f:G(F)\backslash G(\mathbb{A}_F^\infty)/K^p\to L\text{ 连续}\},
\]
它是 $G(F_p)$ 的容许 Banach 空间表示\tn{底本此处作 $G(F_p)$，按上下文应指 $p$ 处各位置对应群的积 $\prod_{v|p}G(F_v)$；译文照排。（第 350 页）}，并且由于我们的群在无穷处紧，它与 $Y(K^p)$ 的完备上同调（基变换到 $L$）一致。该 Banach 空间表示中的局部解析向量由
\begin{equation}\tag{9.6.4}\label{eq:a4-16-9-6-4}
\wh S(K^p,L)^{\mathrm{an}}=\{f:G(F)\backslash G(\mathbb{A}_F^\infty)/K^p\to L\text{ 局部解析}\}
\end{equation}
给出。

如上，记 $N_v\subset B_v$ 为 $B_v$ 的幂幺根，$N_{v,0}=N_v\cap B_{v,0}$。若令 $N_0=\prod\limits_{v|p}N_{v,0}$，则可在 $N_0$-不变量 $\wh S(K^p,L)^{N_0}$ 上定义 $T^+$ 的作用，方式为
\begin{equation}\tag{9.6.5}\label{eq:a4-16-9-6-5}
t.f=\sum_{n\in N_0/tN_0t^{-1}}ntf.
\end{equation}
下面这个命题比较带 $\mathcal A^+(G)$-作用的模 (9.6.3) 与 (9.6.4)（其中 $\mathcal A^+(G_v)$ 经由包含 $L[\Sigma_v^+]\subset L[T_v^+]$ 作用在 (9.6.4) 的 $N_0$-不变量上）。

\begin{prop}{9.6.6}\label{prop:a4-16-9-6-6}
设 $N_0=\prod_vN_{v,0}$。则存在同构
\[
M^\dagger(G,K^p,(\kappa_v)_{v|p})\isom\bigl(\wh S(K^p,L)^{\mathrm{an}}\bigr)^{N_0}[(\kappa_v)_{v|p}],
\]
它对 Atkin--Lehner 环 $\mathcal A^+(G)=\bigotimes_{v|p}\mathcal A^+(G_v)$ 的作用以及远离 $p$ 的 Hecke 作用等变。这里右端的 $(-)[(\kappa_v)_{v|p}]$ 表示 $N_0$-不变量中使 $T^0=\prod_vT_v^0$ 经由 $\prod_v\kappa_v$ 作用的子空间。
\end{prop}

\begin{proof}
例如见 [Loe11, Prop. 3.10.1, Prop. 3.10.2]。
\end{proof}

除了固定的局部解析特征标，还可以把同样的构造用于 $T^0$ 的一族特征标，乃至用于万有特征标
\[
\kappa^{\mathrm{univ}}:T^0\to\Gamma(\wh T^0,\OO_{\wh T^0})^\times,
\]
其中 $\wh T^0$ 表示 $T^0$ 的连续特征标构成的空间。这意味着，在 $M^\dagger(G,K^p,\kappa^{\mathrm{univ}})$ 的定义中，我们把表示 $\Ind_{B_{v,0}}^{I_v}\kappa_v$ 换成 $I_v$-表示 $\Ind_{N_{v,0}}^{I_v}\mathbf 1$（它包含各个特征标 $\kappa_v$ 对应的全部表示 $\Ind_{B_{v,0}}^{I_v}\kappa_v$）。

\begin{prop}{9.6.7}\label{prop:a4-16-9-6-7}
存在同构
\[
M^\dagger(G,K^p,\kappa^{\mathrm{univ}})\isom\bigl(\wh S(K^p,L)^{\mathrm{an}}\bigr)^{N_0},
\]
它对 $L[T^+]$ 的作用以及远离 $p$ 的 Hecke 算子作用等变。
\end{prop}

\begin{proof}
这是 [Fu21, Theorem 1.2] 的特殊情形；另见 [Eme06b, §3.2] 的讨论。
\end{proof}

下面把这些超收敛 $p$-adic 自守形式空间与函子 $\mathfrak A_G^{\mathrm{rig}}$ 猜想性地联系起来。在伽罗瓦一侧，记 $\XX_{G,F,S}=\XX^{\mathrm{rig}}_{G,F,S}$ 为叠 $\XX_{G,F,S}$ 的刚性解析一般纤维。注意 $\XX_{G,F,S}$ 一般只是导出叠（而非经典叠），因此必须调整\ysrc{351}其一般纤维的定义（届时它将是导出刚性解析 Artin 叠），也就是把 (5.1.12) 推广到导出情形。由于我们不打算在此讨论导出刚性解析几何，故不写出 $\XX_{G,F,S}$ 的精确定义。

注意，类似于 (9.2.1)，给定 $v|p$ 时存在态射
\[
f_{\mathrm{rig},v}:\XX_{G,F,S}=(\XX_{G,F,S})^{\mathrm{rig}}_\eta\to\XX_{G,F_v},
\]
它由到 $(\XX_{G,F_v})^{\mathrm{rig}}_\eta$ 的相应映射与态射 (5.1.31) 复合而成。我们记
\[
f_{\mathrm{rig}}:\XX_{G,F,S}\to\prod_v\XX_{G,F_v}
\]
为各 $f_{\mathrm{rig},v}$ 之积。滥用记号，我们也把 $f_{\mathrm{rig}}$ 限制到 $\XX^{\mathrm{odd}}_{G,F,S}=(\XX^{\mathrm{odd}}_{G,F,S})^{\mathrm{rig}}_\eta\subset\XX_{G,F,S}$ 上得到的映射记作 $f_{\mathrm{rig}}$。

记 $\sigma_{\kappa_v}=\Ind_{B_{v,0}}^{I_v}\kappa_v$，$\sigma_v=\Ind_{N_{v,0}}^{I_v}\mathbf 1$。利用这些表示以及函子 $\mathfrak A^{\mathrm{rig}}_{G(F_v)}$，可以定义凝聚层
\begin{equation}\tag{9.6.8}\label{eq:a4-16-9-6-8}
\begin{aligned}
\mathfrak A^{\mathrm{rig}}((\kappa_v)_{v|p})&=\Bigl(\bigotimes_{v\nmid p}\mathfrak A_{G(F_v)}\bigl(c\text{-}\Ind_{K_v}^{G(F_v)}\mathbf 1\bigr)\Bigr)^{\mathrm{rig}}\otimes\bigotimes_{v|p}\mathfrak A^{\mathrm{rig}}_{G(F_v)}\bigl(c\text{-}\Ind_{I_v}^{G(F_v)}\sigma_{\kappa_v}\bigr),\\
\mathfrak A^{\mathrm{univ}}_{\mathrm{rig}}&=\Bigl(\bigotimes_{v\nmid p}\mathfrak A_{G(F_v)}\bigl(c\text{-}\Ind_{K_v}^{G(F_v)}\mathbf 1\bigr)\Bigr)^{\mathrm{rig}}\otimes\bigotimes_{v|p}\mathfrak A^{\mathrm{rig}}_{G(F_v)}\bigl(c\text{-}\Ind_{I_v}^{G(F_v)}\sigma_v\bigr).
\end{aligned}
\end{equation}

\begin{conj}{9.6.9}\label{conj:a4-16-9-6-9}
存在同构
\[
M^\dagger(G,K^p,(\kappa_v)_{v|p})\cong R\Gamma_c\bigl(\XX^{\mathrm{odd}}_{G,F,S},f_{\mathrm{rig}}^*\mathfrak A^{\mathrm{rig}}((\kappa_v)_{v|p}),
\]
\[
M^\dagger(G,K^p,\kappa^{\mathrm{univ}})\cong R\Gamma_c\bigl(\XX^{\mathrm{odd}}_{G,F,S},f_{\mathrm{rig}}^*\mathfrak A^{\mathrm{univ}}_{\mathrm{rig}}\bigr),
\]
\tn{底本第一式末尾缺一个右括号，译文照排；该处应补上右括号。（第 351 页）}它们与 Hecke 作用以及 Atkin--Lehner 环的作用相容。
\end{conj}

\begin{rem}{9.6.10}\label{rem:a4-16-9-6-10}
Atkin--Lehner 环经由 $L[T^+]$ 在 $c\text{-}\Ind_{I_v}^{G_v}\sigma_{\kappa_v}$ 上的典范作用，作用在右端，例如见 [KS12, Lemma 2.2, Lemma 2.3]。

我们指出，这个猜想在某种意义上与猜想 9.3.2 相对偶：我们在对称空间一侧使用上同调（此处对称空间只是点集），而在凝聚层一侧使用紧支上同调（注意叠 $\XX^{\mathrm{odd}}_{G,F,S}$ 通常是某个 Stein 空间对约化群作用的商，因此可以像 [Chi90] 那样、类似于 §7.1 中的用法来定义紧支上同调）。这种对偶性也解释了为何本猜想中使用沿 $f_{\mathrm{rig}}$ 的 $*$-拉回，而不用猜想 9.3.2 中的 $!$-拉回。
\end{rem}

\begin{rem}{9.6.11}\label{rem:a4-16-9-6-11}
在 $M^\dagger(G,K^p,(\kappa_v)_{v|p})$ 的定义中，我们也可以对某些 $v\nmid p$ 的位置插入非平凡的 $K_v$-表示 $\sigma_v$ 作为系数。此时我们预期与猜想 9.6.9 相同的同构仍成立，但 (9.6.8) 中平凡的 $K_v$-表示 $\mathbf 1$ 自然要换成 $\sigma_v$。
\end{rem}

\subsubsection*{9.6.12.\quad 有限斜率空间与 Emerton 局部解析 Jacquet 函子}

我们回顾有限斜率空间以及局部解析 Jacquet 模的定义 [Eme06a]。

空间 $M^\dagger(G,K^p,(\kappa_v)_{v|p})$ 与 $(\wh S(K^p,L)^{\mathrm{an}})^{N_0}$ 都是（紧型的）局部凸拓扑 $L$-向量空间，带有幺半群 $L[T^+]$ 的作用。与局部\ysrc{352}凸 $L$-向量空间 $V$（带有拓扑 $T^+$-作用，或对某个位置 $v|p$ 带有 $T_v^+$-作用）相伴的有限斜率空间 $V_{\mathrm{fs}}$ 的构造，是一种把 $L[T^+]$ 作用延拓为 $T\supset T^+$ 的局部解析表示的典范手续。其形式定义为
\[
V_{\mathrm{fs}}=\Hom_{L[T^+],b}\bigl(\Gamma(\wh T,\OO_{\wh T}),V\bigr),
\]
可对照 [Eme06a, Definition 3.2.1]。这里 $\wh T$ 是 $T$ 的特征标的刚性解析空间（注意 $\Gamma(\wh T,\OO_{\wh T})$ 包含分布代数 $\mathcal D(T)$，但严格更大），下标 $b$ 表示取连续同态构成的空间并赋予强拓扑。于是 $V_{\mathrm{fs}}$ 上的 $\Gamma(\wh T,\OO_{\wh T})$-模结构限制为 $\mathcal D(T)$ 的作用，从而限制为 $T$ 的作用；后者事实上是局部解析的 [Eme06a, Proposition 3.2.4]。若这个表示是（本质）容许的，则有限斜率空间的强对偶是 Frechet--Stein 代数 $\Gamma(\wh T,\OO_{\wh T})$ 上的余容许模。由于 $\wh T$ 是 Stein 空间，这意味着存在（必然唯一的）$\wh T$ 上的凝聚层 $\mathcal M^*_V$，使得
\begin{equation}\tag{9.6.13}\label{eq:a4-16-9-6-13}
(V_{\mathrm{fs}})'\cong\Gamma(\wh T,\mathcal M^*_V)
\end{equation}
作为 $\Gamma(\wh T,\OO_{\wh T})$-模同构。

\begin{rem}{9.6.14}\label{rem:a4-16-9-6-14}
\begin{enumerate}
\item[(i)] 若 $V=\Pi^{N_{v,0}}$，其中 $\Pi$ 是局部解析 $G_v$-表示，则 $V$ 带有由 (9.6.5) 定义的 $T^+$-作用。此时 $V_{\mathrm{fs}}$（作为局部解析 $T_v$-表示）的本质容许性，例如可以通过假设 $\Pi|_H=C^{\mathrm{an}}(H,L)^m$（其中 $H\subset G_v$ 是某个开紧子群，$m>0$）来保证，见 [Eme06a, Proposition 3.2.24] 与 [BHS17b, Proposition 5.3]（后者依赖 [Eme06a]）。在此情形，[BHS17b, Proposition 5.3] 还表明：在 $\wh T_v$ 上局部地，层 $\mathcal M^*_V$ 是 $\wh T^0_v$ 上的有限投射模。
\item[(ii)] 假设 $V_{\mathrm{fs}}$ 上的 $T$-表示是本质容许的。于是对 Stein 空间 $\wh T$ 使用 Serre 对偶（正如对 (7.1.5) 中的开单位圆盘 $\mathbf U$ 所做的那样），可以把 (9.6.13) 改写为
\[
V_{\mathrm{fs}}=R\Gamma_c(\wh T,\mathcal M_V),
\]
其中 $\mathcal M_V=\mathbf D(\mathcal M^*_V)$ 是 $\mathcal M^*_V$ 的（$\wh T$ 上的）Serre 对偶。先验地，这只是凝聚层的复形。不过，在好的情形下（例如 (i) 的情形），层 $\mathcal M^*_V$ 实际上是 Cohen--Macaulay 模，从而 $\mathcal M_V$ 集中在一个次数上（尽管这个次数可能是严格负的，视 $\mathcal M^*_V$ 的支撑的维数而定）。
\end{enumerate}
\end{rem}

\begin{defn}{9.6.15}\label{defn:a4-16-9-6-15}
局部解析 $G_v$-表示 $\Pi_v$ 的局部解析 Jacquet 模 $J_{B_v}(\Pi_v)$，指的是 $(\Pi_v^{N_{v,0}})_{\mathrm{fs}}$ 上的 $T_v$-表示。
\end{defn}

有了这个定义，命题 9.6.7 直接蕴含：

\begin{prop}{9.6.16}\label{prop:a4-16-9-6-16}
存在局部解析 $T$-表示的同构
\[
M^\dagger(G,K^p,\kappa^{\mathrm{univ}})_{\mathrm{fs}}\isom J_B\bigl(\wh S(K^p,L)^{\mathrm{an}}\bigr).
\]
\end{prop}

若只关注 $M^\dagger(G,K^p,\kappa^{\mathrm{univ}})$ 的有限斜率子空间，就可以陈述一个猜想，精确描述这样一个凝聚层：其紧支上同调猜想性地计算 $M^\dagger(G,K^p,\kappa^{\mathrm{univ}})_{\mathrm{fs}}$\ysrc{353}，且不涉及猜想性的函子 $\mathfrak A_G^{\mathrm{rig}}$（当然，这个猜想受这些猜想性函子及其关于抛物诱导的预期性质的启发）。

如上，对 $v\nmid p$ 记 $\mathfrak A_{\mathrm{rig},v}=(\mathfrak A_{G(F_v)}(c\text{-}\Ind_{K_v}^{G_v}\mathbf 1))^{\mathrm{rig}}$。对 $v|p$ 考虑态射
\[
f_{\mathrm{rig},v}\times\iota_v:\XX^{\mathrm{odd}}_{G,F,S}\times\wh T\longrightarrow\XX_{G,F_v}\times\wh T_v,
\]
其中 $\iota_v$ 由到第 $v$ 个因子的投影再复合上 $\wh T_v$ 的自同构
\[
(\delta_1,\dots,\delta_n)\mapsto\delta_{B_v}\cdot\bigl(\delta_1,\delta_2(\varepsilon\circ\mathrm{rec}_{F_v}),\dots,\delta_n(\varepsilon\circ\mathrm{rec}_{F_v})^{n-1}\bigr)
\]
给出，其中 $\delta_{B,v}$ 表示模特征（modulus character）\tn{底本上面的公式中写作 $\delta_{B_v}$，紧接的文字中写作 $\delta_{B,v}$，下标不一致；译文照排。（第 353 页）}。此外考虑如下图表
\[
\begin{tikzcd}
& \overline{\XX}_{B_v}\arrow[dl,"\pi_v"']\arrow[dr,"\alpha_v"] & \\
\XX_{G,F_v}\times\wh T_v & & \XX_{T_v}
\end{tikzcd}
\]
并定义
\[
\mathfrak A_{\mathrm{rig},v}=\pi_{v,*}\alpha_v^*\bigl(\OO_{\XX_{T_v}}([F_v:\Q_p](0,-1,\dots,-n+1))[\mathrm{rk}_{\Q_p}T_v]\bigr).
\]
这里，如同猜想 6.2.4 (2) 中那样，扭 $(0,-1,\dots,-n+1)$ 的含义是：取 $\wh T_v$ 上的层，并通过赋予由代数特征标 $(0,-1,\dots,-n+1)$ 给出的 $\GG_m^n$-等变结构，把它看作 $\XX_{T_v}=\wh T_v/\GG_m^n$ 上的层。

\begin{rem}{9.6.17}\label{rem:a4-16-9-6-17}
注意，谈论 $\pi_{v,!}$ 而非 $\pi_{v,*}$ 会更典范一些。但由定理 5.3.12，$\pi_v$ 是固有的（proper），故这两个函子一致，我们不必操心 $\pi_{v,!}$ 的定义（那多半会涉及凝聚结构（condensed structures））。
\end{rem}

记
\[
g_{\mathrm{rig}}=\prod_{v\nmid p}f_{\mathrm{rig},v}\times\prod_{v|p}g_{\mathrm{rig},v}:\XX^{\mathrm{odd}}_{G,F,S}\times\wh T\to\prod_v\XX_{G,F_v}\times\wh T,
\]
则可陈述如下猜想：

\begin{conj}{9.6.18}\label{conj:a4-16-9-6-18}
存在同构
\[
M^\dagger(G,K^p,\kappa^{\mathrm{univ}})_{\mathrm{fs}}\cong R\Gamma_c\Bigl(\XX_{G,F,S}\times\wh T,g_{\mathrm{rig}}^*\bigl(\bigotimes_v\mathfrak A_{\mathrm{rig},v}\bigr)\Bigr),
\]
它对远离 $p$ 的 Hecke 算子作用、以及 $p$ 处位置的 Atkin--Lehner 环作用等变。
\end{conj}

\begin{rem}{9.6.19}\label{rem:a4-16-9-6-19}
注意 $\mathfrak A_{\mathrm{rig},v}$ 定义中的次数移位是必要的：预期猜想 9.6.18 中的凝聚层的支撑维数为 $\sum_{v|p}\mathrm{rk}_{\Q_p}T_v$。因此有了这个次数移位，就应当预期紧支上同调位于次数 $0$。这也应能解释：为何在 §7.1 的解析情形中我们实际上需要某种次数移位。
\end{rem}

\ysrc{354}
\subsubsection*{9.6.20.\quad $p$-adic 自守形式的 Hecke 特征空间}

我们假设 $\Pi_v$ 是满足注 9.6.14 中假设的 $G_v$-表示，即假设 $\Pi_v|_{H_v}$ 同构于 $C^{\mathrm{an}}(H_v,L)^m$（其中 $m>0$，$H_v\subset G_v$ 是某个紧开子群）。给定 $T$ 的局部解析特征标 $\delta$，可利用按 (9.6.13) 定义的凝聚层 $\mathcal M^*_{\Pi_v}:=\mathcal M^*_{\Pi_v^{N_{v,0}}}$ 计算特征空间
\[
\Hom_T(\delta,J_{B_v}(\Pi_v))=\Hom_{T^+}(\delta,\Pi_v^{N_{v,0}}).
\]
事实上，可以把 $\delta$ 看作 $\wh T$ 中的点，并算出
\begin{equation}\tag{9.6.21}\label{eq:a4-16-9-6-21}
\Hom_T(\delta,J_{B_v}(\Pi_v))=(\mathcal M^*_{\Pi_v}\otimes_{\OO_{\wh T}}k(\delta))',
\end{equation}
它就是 $\mathcal M^*$ 在点 $\delta$ 处的纤维的对偶空间。我们指出，该公式还有一个导出版本：
\[
R\Hom_T(\delta,J_{B_v}(\Pi_v))=\bigl(R\Gamma(\wh T,\mathcal M^*_{\Pi_v}\otimes^{L}_{\OO_{\wh T}}k(\delta))\bigr)'.
\]

\begin{rem}{9.6.22}\label{rem:a4-16-9-6-22}
\begin{enumerate}
\item[(i)] 至少当 $\delta$ 是局部代数特征标时，特征空间 $\Hom_T(\delta,J_{B_v}(\Pi_v))$ 可以用 Breuil 的伴随公式 [Bre15, Remarque 4.4] 来计算：
\[
\Hom_T(\delta,J_{B_v}(\Pi_v))=\Hom_{G_v}\bigl(\mathcal F_{B_v}^{G_v}(\delta),\Pi_v\bigr),
\]
其中 $\mathcal F_{B_v}^{G_v}(-)$ 是 Orlik 与 Strauch 定义的函子，见上文 §6.2.14。粗略地说，这个伴随关系说明：有限斜率部分
\[
M^\dagger(G,K^p,\kappa^{\mathrm{univ}})_{\mathrm{fs}}=J_B\bigl(\wh S(K^p,L)^{\mathrm{an}}\bigr)
\]
由这样的形式 $f\in\wh S(K^p,L)^{\mathrm{an}}$ 组成：它们生成的局部解析子表示的所有子商都出现在抛物诱导表示中。
\item[(ii)] 若 $\delta_v$ 是 $T_v$ 的局部解析特征标，则它通过限制诱导出 $T_v^+$ 的特征标 $\kappa_v$，以及特征标 $\chi:\mathcal A^+(G_v)=L[\Sigma_v^+]\to L$。反过来，任何这样的对 $(\kappa_v,\chi)$ 都典范地定义了 $T_v$ 的特征标 $\delta_v$。用这些记号并利用 (9.6.21)，可以改写为
\begin{align*}
\Hom_T\bigl(\delta,J_B(\wh S(K^p,L)^{\mathrm{an}})\bigr)&=\Hom_{T_v^+}\bigl(\delta,(\wh S(K^p,L)^{\mathrm{an}})^{N_0}\bigr)\\
&=M^\dagger(G,K^p,(\kappa_v)_v)^{\mathcal A^+(G)=\chi}.
\end{align*}
\end{enumerate}
\end{rem}

\subsubsection*{9.6.23.\quad 特征簇}

我们（略微不正式地）说明与特征簇构造之间的关系，后者例如见 [Eme06b]。猜想 9.6.9 中的 Hecke 作用尤其是某个大 Hecke 代数 $\mathbf T$ 的作用，而 $\mathbf T$ 是远离 $p$ 的好位置处的球 Hecke 代数（spherical Hecke algebra）的张量积。记 $\mathbf T(K^p)$ 为 $\mathbf T$ 在
\[
\wh S(K^p,\OO/\varpi^m)=\{f:G(F)\backslash G(\mathbb{A}_F^\infty)/K^p\to\OO/\varpi^m\text{ 局部常值}\}
\]
的自同态环中的像的逆向极限。则 $\mathbf T(K^p)$ 是诺特半局部环，于是可以构造与它相伴的形式概形 $\mathcal Y(K^p)=\Spf\mathbf T(K^p)$，以及它的刚性解析一般纤维 $\mathfrak Y(K^p)=\mathcal Y(K^p)^{\mathrm{rig}}_\eta$。按照 [Eme06b]（或上文注 9.6.14 及其前面的讨论），Jacquet 模 $J_B(\wh S(K^p,L)^{\mathrm{an}})'$ 的对偶自然成为 $\mathfrak Y(K^p)\times\wh T$ 上的凝聚层 $\mathcal M^*(G,K^p)$。这个凝聚层的支撑称为与 $K^p$ 相伴的特征簇 $\mathcal E(G,K^p)$（这里省略了主 Bernstein 块的选取，可对照注 9.6.1）。

\ysrc{355}
\begin{rem}{9.6.24}\label{rem:a4-16-9-6-24}
当然也可以改用 $M^\dagger(G,K^p,\kappa^{\mathrm{univ}})_{\mathrm{fs}}$ 来构造特征簇；这正是 [Che04] 或 [Loe11] 的观点。

如例如 §1.3 所指出的，猜想性地，完备大 Hecke 代数 $\hat{\mathbf T}(K^p)$ 应当与一族伪形变环的乘积一致；这些伪形变环或者被定义为、或者被证明为同构于 $\XX_{G,F,S}$ 底下的经典叠 $\XX^{cl}_{G,F,S}$ 的不变整体截面；也就是说，猜想性地有同构
\[
\hat{\mathbf T}(K^p)\cong\Gamma\bigl(\XX^{cl}_{G,F,S},\OO_{\XX^{cl}_{G,F,S}}\bigr).
\]
记 $\XX^{ps}_{G,F,S}$ 为与 $(G,F,S)$ 对应的各伪形变环的形式谱的不交并。于是，借助映射 $\mathcal Y(K^p)\to\XX^{ps}_{G,F,S}$（它由与自守形式相伴的伽罗瓦表示的构造得到），可以把 $\mathcal M^*(G,K^p)$ 看作 Stein 空间 $\mathfrak X^{ps}_{G,F,S}\times\wh T$ 上的凝聚层，其中 $\mathfrak X^{ps}_{G,F,S}=(\XX^{ps}_{G,F,S})^{\mathrm{rig}}_\eta$。猜想 9.6.18 现在大致是说：凝聚层 $\mathcal M^*(G,K^p)$ 在某种意义上局部化到 $\XX_{G,F,S}\times\wh T$ 上，即该叠上存在一个凝聚层（它是猜想中显式描述的那个凝聚层的 Serre 对偶），其（导出）整体截面集中在次数零，并且与 Stein 空间 $\mathfrak X^{ps}_{G,F,S}\times\wh T$ 上的层 $\mathcal M^*(G,K^p)$ 的整体截面一致。
\end{rem}

\begin{rem}{9.6.25}\label{rem:a4-16-9-6-25}
\begin{enumerate}
\item[(i)] 我们指出，实际上可以把特征空间的计算 (9.6.21) 加强，从而不仅计算 $p$ 处 Hecke 作用（即 Atkin--Lehner 环 $\mathcal A^+(G)$ 的作用）的特征空间，还能计算远离 $p$ 的 Hecke 作用（即 $\hat{\mathbf T}(K^p)$ 中 Hecke 算子的作用）的特征空间。设 $\frakm_\rho\subset\mathbf T(K^p)$ 为与某个 Hecke 特征系相伴的、附有伽罗瓦表示 $\rho:\Gal_{F,S}\to\mathrm{GL}_n(L)$ 的理想。则 $\frakm_\rho$ 定义了 $\mathcal Y(K^p)^{\mathrm{rig}}_\eta$ 的一个点 $x_\rho$，并且发现
\begin{equation}\tag{9.6.26}\label{eq:a4-16-9-6-26}
\begin{aligned}
&\Hom_G\Bigl(\widehat{\bigotimes}_v\mathcal F_B^G(\delta_v),\wh S(K^p,L)^{\mathrm{an}}[\frakm_\rho]\Bigr)\\
&\qquad=\Hom_T\Bigl(\prod_v\delta_v,J_B(\wh S(K^p,L)^{\mathrm{an}})[\frakm_\rho]\Bigr)\\
&\qquad=\bigl(\mathcal M^*(G,K^p)\otimes k(x)\bigr)',
\end{aligned}
\end{equation}
其中 $x=(x_\rho,(\delta_v)_v)\in\mathcal Y(K^p)^{\mathrm{rig}}_\eta\times\wh T$。
\item[(ii)] 在好的情形下，$\mathcal Y(K^p)^{\mathrm{rig}}_\eta$ 与 $\XX_{G,F,S}$ 在过渡到连通分支之后一致（相差一个 $\GG_m$ 的平凡作用），而例如这些分支上的剩余伽罗瓦表示是某个固定的不可约奇表示。在这些情形下（差一个 Serre 对偶），猜想 9.6.18 给出凝聚层 $\mathcal M^*(G,K^p)$ 的一个非常精确的描述，其中只涉及伽罗瓦表示（或 $(\varphi,\Gamma)$-模）的叠。于是我们也得到特征空间 (9.6.26) 的一个（猜想性的）纯伽罗瓦理论的描述。
\end{enumerate}
\end{rem}

\subsubsection*{9.6.27.\quad 更一般的志村簇}

我们回到更一般的情形，即去掉 $G$ 在无穷处紧的假设。特别地，局部对称空间（志村簇）$Y(K_f)$ 可以在多个次数上有上同调，因此与其处理单个上同调群，不如直接使用复形。

\ysrc{356}
我们再次在整除 $p$ 的位置取定分裂极大环面 $T_v$ 与 Borel 子群 $B_v$。设 $\kappa=\prod_{v|p}\kappa_v:T\to L$ 为局部解析特征标。按上面的方式取定 $K^p$ 与 $I_v$，并令 $K_f=K^p\times\prod_{v|p}I_v$。则表示
\[
W(\kappa)=\widehat{\bigotimes}_{v|p}\Ind_{B_{v,0}}^{I_v}\kappa_v
\]
定义了 $Y(K_f)$ 上一个「局部解析」系数系统，仍记作 $W(\kappa)$，见 [Han17, 2.2, 3]（该系数系统 $W(\kappa)$ 在文献中记作 $\mathcal A_\kappa$）。上同调复形 $R\Gamma(Y(K_f),W(\kappa))$ 于是带有 Atkin--Lehner 环 $\mathcal A^+(G)$ 的作用。当然，这个构造还有涉及万有特征标 $\kappa^{\mathrm{univ}}$ 的变体 $W(\kappa^{\mathrm{univ}})$。

\begin{rem}{9.6.28}\label{rem:a4-16-9-6-28}
复形 $R\Gamma(Y(K_f),W(\kappa))$ 应视为固态 $L$-向量空间导出范畴中的对象，特别地它只在拟同构意义下良定义。其上同调群是紧型的 $L$-向量空间。$\mathcal A^+(G)$ 的作用也应视为导出范畴中的作用，即给定元素所定义的自同态只在同伦意义下良定义，等等。
\end{rem}

类似地可以考察完备上同调，关于它在导出意义下的讨论可比较 [Fu21, 5]。我们同样用复形代替单个上同调群，即考虑复形
\[
\wt R\Gamma(Y(K^p),L)=\Bigl(\varprojlim_{m}\varinjlim_{K_p}R\Gamma(Y(K^pK_p),\OO/\varpi^m)\Bigr)\otimes_{\OO}L.
\]
同样把它视为固态 $L$-向量空间导出范畴中的对象。在这个导出范畴中，它带有 $G$ 的作用，且其上同调群是容许 $G$-表示。更好的是，$\wt R\Gamma(Y(K^p),L)$ 的对偶，即计算完备同调的复形，若只把它看作 $\OO[[K]]$-模的复形，则是完美复形。类似于 (9.6.4)，可以过渡到局部解析向量 $\wt R\Gamma(Y(K^p),L)^{\mathrm{an}}$。于是有类似于命题 9.6.6 的比较。

\begin{prop}{9.6.29}\label{prop:a4-16-9-6-29}
存在拟同构
\[
R\Gamma(Y(K_f),W(\kappa^{\mathrm{univ}}))\cong\bigl(\wt R\Gamma(Y(K^p),L)^{\mathrm{an}}\bigr)^{N_0},
\]
它对 Atkin--Lehner 环的作用等变。
\end{prop}

\begin{proof}
这基本上就是 [Fu21, Theorem 1.2]。
\end{proof}

\begin{rem}{9.6.30}\label{rem:a4-16-9-6-30}
当然也有把单个复形 $R\Gamma(Y(K_f),W(\kappa))$（取固定的 $\kappa$ 而非万有特征标）作比较的版本。作这种比较时，右端须取导出的特征空间。这里略去技术细节。
\end{rem}

Fu [Fu21, 6] 延续了与 §9.6.2 的平行之处，并基于复形 $(\wt R\Gamma(Y(K^p),L)^{\mathrm{an}})^{N_0}$（或其对偶）给出特征簇的导出构造。我们在此不讨论这一点，但要指出：（至少在启发式层面）有限斜率部分 $(\wt R\Gamma(Y(K^p),L)^{\mathrm{an}})^{N_0}_{\mathrm{fs}}$ 的对偶，猜想性地应当局部化为 $\XX_{G,F,S}\times\wh T$ 上凝聚层导出范畴中的对象\tn{底本此处把该叠记作 $\mathfrak X_{G,F,S}$（花体 X），与 §9.6.2 中的 $\mathcal X_{G,F,S}$ 字体不一致；译文按同一对象统一记作 $\mathcal X_{G,F,S}$。（第 356 页）}。与其把特征簇的构造说得更精确，不如在此背景下陈述下述与猜想 9.6.9（从而与猜想 9.3.2）类似的猜想。与 §9.6.2 类似，我们有叠 $\XX_{G,F,S}$ 以及整体到局部的映射 $f_{\mathrm{rig}}$。

\ysrc{357}
\begin{conj}{9.6.31}\label{conj:a4-16-9-6-31}
设 $d$ 为 $Y(K_f)$ 的维数。则存在同构
\[
R\Gamma(Y(K_f),W(\kappa))[d]\cong R\Gamma_c\bigl(\XX^{\mathrm{odd}}_{G,F,S},f_{\mathrm{rig}}^*\mathfrak A^{\mathrm{rig}}((\kappa_v)_{v|p})\otimes V_\mu\bigr),
\]
它对（导出的）Hecke 作用以及 Atkin--Lehner 环的作用相容。
\end{conj}

\begin{rem}{9.6.32}\label{rem:a4-16-9-6-32}
事实上，如同 §9.3.1 那样，应当把 $G$ 换成它到 $\Q$ 的 Weil 限制，从而上同调复形上还带有 $\Gal_F$ 的额外作用（此时 $F$ 是志村数据的反射域）。与猜想 9.3.2 一样，猜想 9.6.31 中的同构届时也应对 $\Gal_F$-作用等变。
\end{rem}

\subsubsection*{9.6.33.\quad 伴点与 Breuil 基座猜想}

我们回到 §9.6.2 的设定，并额外假设 $G$ 是定酉群（当然，在更一般的背景下预期有类似结果）。在 [Bre15, Conj. 5.3, Conj. 6.1] 中，Breuil 对（容许的）局部解析表示 $\hat S(G,K^p)^{\mathrm{an}}[\frakm_\rho]$ 的基座（即最大的半单子表示）给出了精确猜想，其中伽罗瓦表示 $\rho$（系数在 $L$ 中）在整除 $p$ 的位置为晶体，并且 Weil--Deligne 表示 $\mathrm{WD}(\rho|_{\Gal_{F_v}})$ 上晶体 Frobenius 的特征值 $\varphi_{1,v},\dots,\varphi_{n,v}\in L$ 满足：当 $i\neq j$ 时 $\varphi_i/\varphi_j\notin\{1,q_v\}$。Breuil 的猜想描述了局部解析主级数表示 $(\Ind_{B_v}^{G(F_v)}\delta_v)^{\mathrm{an}}$ 的哪些不可约子商（按整除 $p$ 的位置作张量积）嵌入空间 $\hat S(G,K^p)^{\mathrm{an}}[\frakm_\rho]$。

该猜想的一个弱化版本，是精确描述满足
\[
(\rho,(\delta_v)_v)\in\mathcal E(G,K^p)\subset\mathfrak X^{ps}_{G,F,S}\times\prod_v\wh T_v
\]
的特征标 $(\delta_v)_v\in\prod_v\wh T_v$。我们给出这个弱化猜想的精确版本，而 Breuil 基座猜想的精确表述见 [Bre15]。利用注 9.6.25 容易证明：Breuil 的猜想蕴含此处给出的弱化版本。

我们继续假设 $\rho_v=\rho|_{\Gal_{F_v}}$ 是晶体的，且 $\mathrm{WD}(\rho_v)$ 上 Frobenius 特征值两两不同并满足上述正则性假设。假设 $\rho_v$ 有正则 Hodge--Tate 权 $\lambda_v$，并记 $\xi_v$ 为相应的最高权。对每个 Frobenius 稳定旗 $\mathcal F_v\subset\mathrm{WD}(\rho_v)$（定义在 $\overline{\Q}_p$ 上），可以联系
\begin{itemize}
\item 一个非分歧特征标
\[
\delta_{\mathcal F_v}=\mathrm{unr}_{\varphi_{1,v}}\otimes\cdots\otimes\mathrm{unr}_{\varphi_{n,v}}:T_v\to L^\times,
\]
其中 $\varphi_{1,v},\dots,\varphi_{n,v}$ 是由 $\mathcal F_v$ 给出的 Frobenius 特征值的排序；
\item 对每个嵌入 $\tau:F_v\hookrightarrow\overline{\Q}_p$，一个元素 $w_\tau\in\mathcal S_n$，它给出 $D_{\mathrm{dR}}(\rho_v)\otimes_{F,\tau}\overline{\Q}_p$ 上的 Hodge 滤过（的 $\tau$-部分）相对于 $\mathcal F_v$ 的相对位置。记 $w(\mathcal F_v)$ 为元组 $(w_\tau)_\tau\in\prod_{\tau:F_v\hookrightarrow\overline{\Q}_p}\mathcal S_n$。
\end{itemize}

\ysrc{358}
\begin{conj}{9.6.34}\label{conj:a4-16-9-6-34}
设 $\rho$ 是如上所述的表示，且 $(\delta_v)_{v|p}\in\wh T(L)$。则 $(\rho,(\delta_v)_v)\in\mathcal E(G,K^p)$ 当且仅当每个 $\delta_v$ 都具有形式
\[
\delta_v=z^{w_vw_0\cdot\xi_v}\delta_{\mathrm{sm}},
\]
其中 $(\mathcal F_v)_{v|p}$ 是某一族 Frobenius 稳定旗的选取，外尔群元素 $w_v=(w_\tau)_{\tau:F_v\hookrightarrow L}$ 是相应的选取，且满足 $w(\mathcal F_v)\preceq w_v$。这里 $w_0$ 表示外尔群的最长元素，$\preceq$ 表示通常的 Bruhat 序。
\end{conj}

我们指出，这个猜想会是猜想 9.6.18 的直接推论：只需计算猜想 9.6.18 右端出现的凝聚层的支撑。然而，这个支撑基本上可以用定理 5.3.19 算出。

\begin{rem}{9.6.35}\label{rem:a4-16-9-6-35}
对猜想中的 $w_v$，点
\[
x_{w_v}=(\rho,(\delta_{\mathcal F_v}z^{w_vw_0\cdot\xi_v})_v)
\]
通常称为点 $x_{w_0}=(\rho,(\delta_{\mathcal F_v}z^{\xi_v})_v)$ 的\emph{伴点}（companion points）。注意后一个点猜想性地对应于某个经典自守形式 $f$ 的（Hecke 特征系）。由 $x_{w_v}$ 所定义的 Hecke 特征系中的 $p$-adic 自守形式，则通常称为 $f$ 的\emph{伴形式}（companion forms）。

在附加假设下，即在与 3 中的拼接构造相关的所谓 Taylor--Wiles 假设\tn{底本此处作「in 3」，疑指第 3 节；译文照排原文。（第 358 页）}下，该猜想已在 [BHS19, Theorem 5.3.3, Remark 5.3.2] 中得到证明。
\end{rem}

\begin{thm}{9.6.36}\label{thm:a4-16-9-6-36}
假设
\begin{itemize}
\item $p>2$；
\item $F$ 的 CM 扩张 $E$ 非分歧，且不含 $p$ 次单位根；
\item 群 $G$ 在有限位置拟分裂；若 $v$ 是 $F$ 中在 $E$ 里惰性的位置，则 $K_v$ 超特殊；
\item $\rho$ 的剩余伽罗瓦表示 $\bar\rho$ 是 adequate（充分）的。
\end{itemize}
则猜想 9.6.34 成立。
\end{thm}

\begin{rem}{9.6.37}\label{rem:a4-16-9-6-37}
\begin{enumerate}
\item[(i)] 事实上，猜想 6.2.4 (iii) 与猜想 9.6.9 蕴含局部解析基座情形的完整猜想；而 [BHS19, Theorem 5.3.3] 在上述定理的假设下证明了完整的基座猜想。其主要输入是局部解析 Breuil--Mézard 猜想的（证明），即定理 6.2.37。
\item[(ii)] 在非正则 Hodge--Tate 权的情形，上述定理有 Zhixiang Wu [Wu24] 给出的推广；对更一般 Bernstein 分支相伴的特征簇，则有 Breuil--Ding [BD24] 的推广。
\end{enumerate}
\end{rem}

\subsubsection*{9.6.38.\quad 经典与非经典 $p$-adic 自守形式}

如同 §9.6.33，我们继续在定酉群 $G$ 的固定整体设定下工作。设 $x=(\rho,(\delta_v)_v)\in\mathfrak X^{ps}_{G,F,S}\times\prod_v\wh T_v$ 是给定点。我们假设特征标 $\delta_v=z^{\xi_v}\delta_{v,\mathrm{sm}}$ 是局部代数的，其光滑部分 $\delta_{v,\mathrm{sm}}$ 非分歧。回顾注 9.6.22 与注 9.6.25 可知，可以把 $\delta$ 看作由特征标 $\kappa=(\kappa_v)_v:T^0=\prod_{v|p}T_v^0\to L^\times$ 与特征标
\[
\chi=\bigotimes_v\chi_{\delta_{v,\mathrm{sm}}}:\mathcal A^+(G)\to L
\]
组成的二元组 $((\kappa_v)_v,\chi)$。

% 本片止于印刷页 358 页末（§9.6.38 中途），续片自 PDF 372 页起
